\section{Implementation}

The implementation of the real-time Mindfulness Index (MI) system centers on the integration of advanced signal processing with an immersive, adaptive virtual reality (VR) environment. A key innovation of this project is the use of a VR headset and panoramic projection to create a multisensory space, where both visual and auditory feedback are dynamically modulated in response to the user’s physiological state. The sound environment, including spatialized audio and natural soundscapes, is tightly coupled with the projected visuals to reinforce mindfulness and emotional regulation. This holistic approach ensures that the user is fully engaged, with the environment adapting in real time to support the experimental protocol and maximize the effectiveness of neurofeedback.
\subsection{Real-time Production}
\label{sec:realtime_production}

The core objective of the system is to enable real-time estimation and feedback of mindfulness and affective states using EEG and EDA signals. Real-time production refers to the continuous, low-latency processing of biosignals, feature extraction, and computation of indices that reflect the user’s cognitive and emotional state.

\subsection*{Design of the Real-time Experience}
The real-time system is implemented as an immersive Virtual Reality (VR) experience. The user is equipped with a VR headset, an EEG cap, and an EDA sensor, and is situated in a sonorized space designed to minimize external distractions. The environment features a 270-degree panoramic video projection that simulates a natural, ethereal landscape, drawing on principles of biophilic design and multisensory engagement as described in the project documentation (see Section 4.5 and 4.5.5 of the reference). The visual system is complemented by spatialized audio and natural soundscapes, further enhancing the sense of presence and immersion. This setup is intended to facilitate deep states of mindfulness and emotional regulation, leveraging the restorative effects of nature-inspired environments and multisensory integration (see also Section 4.5.7 and 4.5.8).

During the session, the user’s physiological signals are continuously monitored and processed. The VR environment can adapt in real time to the user’s affective state, for example by modulating visual or auditory elements in response to changes in the Mindfulness Index (MI) or Emotional Mindfulness Index (EMI). This closed-loop feedback is designed to support self-regulation and enhance the efficacy of mindfulness training in an engaging, user-centered manner.

\subsection{Key Concepts}
\begin{itemize}
    \item \textbf{Real-time Feedback:} The system provides immediate feedback to the user by continuously updating the Mindfulness Index (MI), Emotional Mindfulness Index (EMI), and Attention Index (ATT) every second.
    \item \textbf{Windowed Processing:} Data is processed in 1-second windows (250 samples at 250 Hz), balancing temporal resolution and noise reduction.
    \item \textbf{Dual Calibration:} Each session begins with two calibration phases (relaxed and focused), allowing the system to adapt to individual baselines and maximize sensitivity to state changes.
    \item \textbf{Continuous Output:} The indices are streamed via LSL, enabling integration with other applications or real-time visualization tools.
\end{itemize}

\begin{table}[h!]
\centering
\caption{Real-time Indices and Their Interpretation}
\begin{tabular}{|l|l|c|l|}
\hline
\textbf{Index Name} & \textbf{Description} & \textbf{Range} & \textbf{Interpretation} \\
\hline
Mindfulness Index (MI) & Personalized measure of mindfulness, combining EEG and EDA features & 0.1--0.9 & Higher values indicate greater mindfulness \\
Emotional Mindfulness Index (EMI) & Focuses on emotional regulation and arousal & 0.05--0.95 & Higher values indicate emotional balance \\
Attention Index (ATT) & Reflects sustained attention (theta power) & 0.05--0.95 & Higher values indicate greater attention \\
\hline
\end{tabular}
\end{table}

\subsection{Signal Processing \\& Communication}
\label{sec:signal_processing}

Signal processing is central to the system’s ability to extract meaningful features from raw biosignals. The following concepts and methods are emphasized:

\subsection{Key Concepts}
\begin{itemize}
    \item \textbf{Artifact Rejection:} Median filtering is applied to both EEG and EDA signals to suppress transient artifacts and noise.
    \item \textbf{Band Power Extraction:} Welch’s method is used to compute power in specific frequency bands relevant to mindfulness and attention.
    \item \textbf{Feature Normalization:} Features are normalized using empirically derived minimum and maximum values, ensuring comparability across sessions and users.
    \item \textbf{Feature Set:} The system extracts nine features per window, each mapped to a cognitive or affective construct.
\end{itemize}

\begin{table}[h!]
\centering
\caption{Extracted Features and Their Ranges}
\begin{tabular}{|l|l|l|c|c|l|}
\hline
\textbf{Feature Name} & \textbf{Channel(s)} & \textbf{Frequency Band (Hz)} & \textbf{Min} & \textbf{Max} & \textbf{Role} \\
\hline
theta\_fz & Fz & 4--8 & 21 & 257,435 & Attention regulation \\
beta\_fz & Fz & 13--30 & 20 & 114,651 & Effortful control \\
alpha\_c3 & C3 & 8--13 & 22 & 1,319,565 & Left body awareness \\
alpha\_c4 & C4 & 8--13 & 27 & 2,554,539 & Right body awareness \\
faa\_c3c4 & C3, C4 & 8--13 (asymmetry) & 0.15 & 1.42 & Emotion regulation (asymmetry) \\
alpha\_pz & Pz & 8--13 & 29 & 3,823,299 & DMN suppression \\
alpha\_po & PO7, PO8 & 8--13 (average) & 29 & 2,210,734 & Visual detachment \\
alpha\_oz & Oz & 8--13 & 15 & 1,797,776 & Relaxation \\
eda\_norm & EDA & --- & 1.972 & 2.005 & Arousal/stress (normalized EDA) \\
\hline
\end{tabular}
\end{table}

\subsection*{Communication}
\begin{itemize}
    \item \textbf{Lab Streaming Layer (LSL):} All indices are transmitted in real time using LSL, facilitating interoperability with other research tools and real-time visualization.
    \item \textbf{Session Logging:} All computed indices are logged for post-hoc analysis and validation.
\end{itemize}

\subsection{Affective State Classification}
\label{sec:affective_classification}

Affective state classification is achieved through a combination of feature engineering, normalization, and adaptive thresholding. The system is designed to map physiological features to interpretable indices that reflect the user’s affective and cognitive state.

\subsection*{Key Concepts}
\begin{itemize}
    \item \textbf{Personalized Thresholds:} Dual calibration establishes user-specific baselines for relaxed and focused states, ensuring that MI and related indices are sensitive to individual differences.
    \item \textbf{Weighted Feature Combination:} Each feature is assigned a weight reflecting its theoretical and empirical relevance to mindfulness and affective regulation.
    \item \textbf{Sigmoid Transformation:} The weighted sum of normalized features is passed through a sigmoid function to produce the final MI, ensuring bounded, interpretable output.
\end{itemize}

\begin{table}[h!]
\centering
\caption{Feature Weights for Mindfulness Index Calculation}
\begin{tabular}{|l|c|l|}
\hline
\textbf{Feature} & \textbf{Weight} & \textbf{Rationale} \\
\hline
theta\_fz & 0.20 & Attention regulation \\
beta\_fz & 0.10 & Effortful control \\
alpha\_c3 & 0.12 & Body awareness (left) \\
alpha\_c4 & 0.12 & Body awareness (right) \\
faa\_c3c4 & 0.25 & Emotion regulation (asymmetry) \\
alpha\_pz & 0.08 & DMN suppression \\
alpha\_po & 0.08 & Visual detachment \\
alpha\_oz & 0.03 & Relaxation \\
eda\_norm & -0.15 & Arousal/stress (negative contribution) \\
\hline
\end{tabular}
\end{table}

\subsection*{Calculation Formula}
The Mindfulness Index (MI) is computed as follows:

\begin{enumerate}
    \item \textbf{Normalization:}
    \[
        \text{norm}_i = \frac{\text{feature}_i - \text{min}_i}{\text{max}_i - \text{min}_i}
    \]
    (clipped to [0, 1] for each feature)
    \item \textbf{Weighted Sum:}
    \[
        S = \sum_{i=1}^9 w_i \cdot \text{norm}_i
    \]
    \item \textbf{Sigmoid Transformation:}
    \[
        \text{MI} = \frac{1}{1 + e^{-4(S - 0.5)}}
    \]
    (clipped to [0.1, 0.9] to avoid saturation)
\end{enumerate}

\begin{table}[h!]
\centering
\caption{Example Calculation (Hypothetical Feature Values)}
\begin{tabular}{|l|c|c|c|c|}
\hline
\textbf{Feature} & \textbf{Value} & \textbf{Normalized} & \textbf{Weight} & \textbf{Contribution} \\
\hline
theta\_fz & 100,000 & 0.386 & 0.20 & 0.077 \\
beta\_fz & 50,000 & 0.436 & 0.10 & 0.044 \\
alpha\_c3 & 500,000 & 0.378 & 0.12 & 0.045 \\
alpha\_c4 & 1,000,000 & 0.391 & 0.12 & 0.047 \\
faa\_c3c4 & 0.8 & 0.561 & 0.25 & 0.140 \\
alpha\_pz & 1,000,000 & 0.262 & 0.08 & 0.021 \\
alpha\_po & 1,000,000 & 0.453 & 0.08 & 0.036 \\
alpha\_oz & 500,000 & 0.269 & 0.03 & 0.008 \\
eda\_norm & 2.0 & 0.848 & -0.15 & -0.127 \\
\hline
\textbf{Total S} & & & & 0.291 \\
\hline
\end{tabular}
\end{table}

Applying the sigmoid:
\[
    \text{MI} = \frac{1}{1 + e^{-4(0.291 - 0.5)}} \approx 0.35
\]

\subsection*{Conceptual Highlights}
\begin{itemize}
    \item \textbf{Affective State as a Continuum:} The MI and EMI indices reflect a continuum of mindfulness and emotional regulation, rather than binary states.
    \item \textbf{Physiological Basis:} Each feature is grounded in neuroscientific and psychophysiological research, linking specific EEG/EDA patterns to cognitive and affective processes.
    \item \textbf{Personalization:} Adaptive thresholds ensure that feedback is meaningful for each user, accounting for individual variability in baseline physiology.
\end{itemize}

% Mindfulness Index Calculation Approaches

In this project, two main approaches for calculating the Mindfulness Index (MI) were implemented and compared:

\subsection*{A. Weighted Feature Sum with Sigmoid Normalization (Current Real-time System)}
The real-time system uses a weighted sum of nine normalized features (EEG and EDA), followed by a sigmoid transformation to ensure the MI remains in the [0,1] range:

\begin{align}
    S &= \sum_{i=1}^9 w_i \cdot \text{norm}_i \\
    \text{MI} &= \frac{1}{1 + e^{-4(S - 0.5)}}
\end{align}

where $\text{norm}_i$ is the normalized value of feature $i$:

\begin{equation}
    \text{norm}_i = \frac{\text{feature}_i - \text{min}_i}{\text{max}_i - \text{min}_i}
\end{equation}

This approach allows for flexible weighting and robust normalization based on empirical feature ranges, supporting real-time adaptation and closed-loop feedback.

\subsection*{B. Classic Pipeline: Linear Combination with Logistic Normalization}
The classic pipeline (see \texttt{eeg\_mindfulness\_index.py}) uses a linear combination of five main features, with both positive and negative weights, followed by a logistic normalization:

\begin{align}
    \text{MI}_{\text{raw}} &= w_1 \cdot \Theta_{Fz} + w_2 \cdot \Alpha_{PO} + w_3 \cdot \text{FAA} - w_4 \cdot \text{Beta}_{\text{Frontal}} - w_5 \cdot \text{EDA}_{\text{norm}} \\
    \text{MI} &= \frac{1}{1 + \exp(-\text{MI}_{\text{raw}} + 1)}
\end{align}

with default weights:

\begin{align*}
    w_1 &= 0.25 \quad \text{(Frontal Theta)} \\
    w_2 &= 0.25 \quad \text{(Posterior Alpha)} \\
    w_3 &= 0.20 \quad \text{(Frontal Alpha Asymmetry)} \\
    w_4 &= 0.15 \quad \text{(Frontal Beta)} \\
    w_5 &= 0.10 \quad \text{(EDA)}
\end{align*}

This formula emphasizes a smaller set of features and uses a simpler normalization, suitable for batch analysis and offline reporting.

\subsection*{Comparison and Rationale}
The real-time system's approach provides greater flexibility and personalization by:
\begin{itemize}
    \item Incorporating a broader set of features (9 vs. 5), including additional EEG bands and spatial locations.
    \item Using empirically derived normalization for each feature, improving robustness across sessions and users.
    \item Allowing for adaptive weighting and dynamic thresholding, supporting closed-loop feedback in immersive environments.
\end{itemize}

The classic pipeline remains valuable for its interpretability and simplicity, especially in retrospective data analysis and model validation.
